\section{Step 8 Recognize the two-byte records}\label{vm-s08}

Select size two for \texttt{PUSH8}, \texttt{JMP} and \texttt{JZ}. Their second byte is payload. A branch payload will later be interpreted as an unsigned offset, while a short push payload is a signed value. Both interpretations share the same structural requirement that one operand byte exists.

\Needspace{20\baselineskip}
\subsection{Write or inspect}
Replace the temporary branch in \texttt{format\_short}. Preserve surrounding labels and the provided frame and exits.

\Needspace{13\baselineskip}
\MARSCodeLabel{MIPS assembly}
\begin{minted}[samepage=true]{mars-idea}
    li $t1, 2
    li $t2, 1
    beq $t9, $t2, format_size
    nop
    li $t2, 48
    beq $t9, $t2, format_size
    nop
    li $t2, 49
    beq $t9, $t2, format_size
    nop
\end{minted}
\subsection{Expected observation}
\texttt{PUSH8} followed by seven selects encoded size two.

\Needspace{11\baselineskip}
\subsection{Check your understanding}
Which statement follows from selecting size two?

\begin{itemize}
  \item \textbf{A} One operand byte must still pass the remaining-length check.
  \item \textbf{B} Its operand is necessarily a valid branch target.
  \item \textbf{C} The guest instruction has already retired.
\end{itemize}
